%!TEX root = ../csuthesis_main.tex
% 设置中文摘要
\keywordscn{Chrono仿真算法； CarlaUE4仿真平台； 多体物理仿真模型}
%\categorycn{TP391}
\begin{abstractzh}
	
	随着多体系统仿真的技术越来越广泛的应用到智慧交通、机械工程等各种各样的领域当中，对于一些要求非常高精度以及高效实时度的问题也变得越发突出起来。在现实中的实际工程应用当中，传统的多体物理仿真的模型有着不能完全还原场景的情况出现、求解的动力学精确程度不高、对多个体之间关系的处理不够到位、无法和虚拟环境建立起良好的联动等方面的问题，很难满足对于复杂多体系统的仿真要求。
	
	本文主要根据 CarlaUE4 仿真平台、Chrono 仿真算法以及 PyCharm 软件来进行实验的设计和模型的创建，并且为了得到高精度并且可靠的多体物理仿真模型，对多体物理仿真模型进行了搭建并完成了 Chrono 算法的优化和平台耦合调试之后，本文便开展了优化前后的对比实验来考察优化后的仿真实现的效果——分析优化前后模型的多体运动轨迹、动力学响应速度等方面的差别，从而证实经过优化以后所建立的仿真模型可以提高很多方面的准确性与真实性，符合复杂的多体环境应用的需求。
	
	最后本文对在CarlaUE4和Chrono基础上所构建的多体物理仿真模型进行系统的建模、参数调整等，总结得到研究结果并指出目前的研究中存在的问题，在此基础上展望未来研究生职业生涯的发展方向，致力于通过不断的积累来提高多体物理仿真的精度、复杂约束的处理以及跨平台的协同仿真技术，并探索它在其他工程领域中的应用前景。
	
	
	
	
	
	
\end{abstractzh}